\documentclass[]{fairmeta}
% Option "twocolumn" available, but please prioritize single-column

\usepackage{amsmath}
\usepackage{amssymb}
\usepackage{amsthm}

% 
\usepackage{algorithm}
\usepackage{algorithmic}
\usepackage{sidecap}


\newcommand{\dayal}[1]{\textcolor{teal}{[Dayal: #1]}}
\newcommand{\mike}[1]{\textcolor{red}{[Mike: #1]}}
\newcommand{\andrey}[1]{\textcolor{red}{[Andrey: #1]}}

\theoremstyle{plain}
\newtheorem{theorem}{Theorem}[section]
\newtheorem{lemma}[theorem]{Lemma}
\newtheorem{proposition}[theorem]{Proposition}
\newtheorem{result}[theorem]{Result}
\newtheorem{corollary}[theorem]{Corollary}

\theoremstyle{definition}
\newtheorem{definition}{Definition}[section]

\newtheorem{property}{Property}[section]

\newtheorem{assumption}{Assumption}[section]

\newtheorem{desideratum}{Desideratum}

% Configure cleveref names
\crefname{definition}{definition}{definitions}
\Crefname{definition}{Definition}{Definitions}


\crefname{property}{property}{properties}
\Crefname{property}{Property}{Properties}

\crefname{assumption}{assumption}{assumptions}
\Crefname{assumption}{Assumption}{Assumptions}
\Crefname{proposition}{Proposition}{Propositions}
\Crefname{result}{Result}{Results}


\crefname{desideratum}{desideratum}{desiderata}
\Crefname{desideratum}{Desideratum}{Desiderata}





\title{A Scalable Measure of Loss Landscape Curvature for Analyzing the Training Dynamics of LLMs}
% Previous title: A Scalable Measure of Loss Landscape Curvature for Analyzing the Training Dynamics of LLMs

% I am hesitent to say a 'novel'

\author[1,2,*]{Dayal Singh Kalra}
\author[1]{Jean-Christophe Gagnon-Audet}
\author[1]{Andrey Gromov}
\author[1]{Ishita Mediratta}
\author[1]{Kelvin Niu}
\author[1]{Alexander H Miller}
\author[1]{Michael Shvartsman}

\affiliation[1]{FAIR at Meta}
\affiliation[2]{University of Maryland, College Park}

% \contribution[*]{Work done as an intern at Meta}
% \contribution[\dagger]{Joint last author}

\abstract{
Understanding the curvature evolution of the loss landscape is fundamental to analyzing the training dynamics of neural networks. 
The most commonly studied curvature measure is Hessian sharpness  ($\lambda_{\max}^H$), defined as the largest eigenvalue of the loss Hessian, which interacts with the learning rate throughout training. 
Despite its significance in analyzing the training dynamics of neural networks, direct measurement of Hessian sharpness remains prohibitive due to extensive computational cost, especially for large-scale models. 
%
To address this challenge, we propose and analyze a computationally efficient measure, \emph{critical sharpness} ($\lambda_c$), which is easily scalable to large models, requiring only about 5–6 forward passes, given the update direction $\Delta \bm{\theta}$. 
Importantly, critical sharpness exhibits several well documented Hessian sharpness phenomena, such as progressive sharpening and Edge of Stability.
We demonstrate the utility of critical sharpness by providing the first demonstration of these sharpness phenomena at unprecedented scale, up to 7$B$ parameters, spanning pre-training and mid-training, and analyze how sharpness dynamics relates to training success. 
More broadly, our results highlight the importance of understanding curvature measures for developing more robust and effective optimization strategies in large-scale neural network training.
}

\date{\today}
\correspondence{\email{dayal@umd.edu}}

% You can add additional metadata fields as follows 
% \metadata[Code]{\url{https://github.com/facebookresearch/repo}}
% \metadata[Blogpost]{\url{https://ai.meta.com/blog/?page=1}}

\begin{document}

\maketitle

\section{Introduction}
\label{section:introduction}

Understanding the evolution of the loss landscape $L(\bm{\theta})$ is fundamental to analyzing the training dynamics of neural networks. The loss landscape describes how the objective function changes with model parameters, and its geometry directly influences optimization, generalization and stability throughout training~\citep{gilmer2022a,kalra2024why}. Intuitively, gradient-based optimization methods can swiftly navigate to the minima in smooth landscapes, while for rough landscapes, such methods can get trapped, hindering convergence or potentially converging to suboptimal solutions~\citep{losslandscapes}. 


Loss landscapes are commonly examined through the eigenvalues and eigenvectors of its Hessian matrix $H(\bm{\theta}):=\nabla_\theta^2 L(\bm{\theta})$. These eigenvalues provide insights into the local curvature of the loss landscape \textemdash large eigenvalues correspond to sharp, steep directions, while small eigenvalues indicate flat, smooth regions. 
The largest eigenvalue of the Hessian, which is often referred to as \emph{Hessian sharpness} $\lambda^H_{\max}$, quantifies the worst-case curvature of the landscape and is a fundamental metric in optimization. Its reciprocal, \emph{flatness} $1/\lambda_{\max}^H$, is a complementary measure used to describe the curvature.
%In convex optimization with vanilla Gradient Descent (GD), the maximum trainable learning rate is determined by the sharpness through $\eta_{\max} < 2 /\lambda_{\max}$. 
In neural network optimization, Hessian sharpness relates to the local training stability~\citep{Wu2018_how_sgd,lewkowycz2020largelearningratephase}. For vanilla Gradient Descent (GD), if the learning rate exceeds the threshold $\sim 2/{\lambda_{\max}^H}$, loss increases as training `catapults' out of the local basin and eventually converges to a flatter region that can accommodate the learning rate~\citep{lewkowycz2020largelearningratephase}. 
% 1. Talk about different optimizers?
This observation generalizes to more complex optimization problems, including mini-batch settings and adaptive optimizers, albeit with different notions of sharpness~\citep{Wu2018_how_sgd,agarwala2025highdimensionalanalysisreveals,cohen2024adaptivegradientmethodsedge,kalra2024why}. 

% 2. Talk about different sharpness phenomena
Hessian sharpness by itself exhibits several robust trends throughout neural network training, particularly documented under constant learning rates. 
% For constant learning rate schedules
Training typically begins with an early reduction in Hessian sharpness~\citep{kalra2023phase,kalra2025universal} followed by a continuous increase until it reaches the stability threshold~\citep{Jastrzebski2020The,cohen2021gradient}. 
Once this threshold is reached, the training stabilizes through a self-stabilization mechanism~\citep{damian2023selfstabilization,cohen2025understanding} \textemdash instead of diverging as classical optimization would predict, Hessian sharpness begins to oscillates around this critical value. The continual increase in Hessian sharpness is referred to as \emph{progressive sharpening} while the subsequent oscillations around the critical threshold is termed the \emph{Edge of Stability}~\citep{cohen2021gradient}.
% 3. relation to learning rate schedules
For the more realistic setting of learning rate schedules involving warmup and decay, sharpness closely follows the learning rate schedule~\citep{gilmer2022a,kalra2023phase,cohen2021gradient}. Due to its particularly close relationship with learning rate, sharpness can used as a diagnostic tool for identifying training instabilities, especially the ones detrimental for training.


% 4. Relation to generalization and sampling ratios?
Beyond training dynamics, sharpness of the final solutions has been related to generalization properties, motivated by the intuition that flatter minima generalize better~\citep{Hochreiter1997FlatM}. However, this relationship has been called into question~\citep{pmlr-v187-kaur23a}, as empirical analysis showing mixed results \textemdash
while flatter solutions generalize better in some settings, they hurt performance in others. 
% 5. Relation to initialization?
Sharpness also plays a role at the beginning of training, with flatter initializations achieving better performance~\citep{NEURIPS2019_876e8108,kalra2024why}.

Despite its value for examining model initialization, training dynamics and generalization properties, analyzing sharpness at scale remains challenging due to extensive computational and memory costs. 
For a model with $n$ parameters, computing the Hessian sharpness requires $\Theta(n^2)$ computational complexity, making such analysis feasible only at small scales or datasets.
As a result, most existing studies are restricted to small-scale experiments, leaving open questions about how sharpness evolves in Large Language Models (LLMs) at scale, and how it relates to optimization and downstream performance.
In this work, we address this challenge by leverage the relationship between curvature and training instability to introduce a scalable measure of loss landscape curvature for analyzing the LLM training dynamics.

\subsection{Contributions}

\begin{itemize}
    \item We propose and analyze critical sharpness, defined as $\lambda_c = 2 / \eta_c$, where $\eta_c$ is the \emph{critical learning rate}\textemdash the smallest learning rate that causes the training loss to increases in the next training step. To estimate the critical learning rate $\eta_c$, we perform an efficient line search along the update direction $\Delta \bm{\theta}$ from training, which utilizes the same computational tools as  large-scale distributed training.
    This simple line search scales linearly with the number of model parameters $\Theta(n)$, in contrast to the quadratic scaling $\Theta(n^2)$ required for computing Hessian sharpness. In practice, we find that this procedure reliably estimates critical sharpness in only $5$-$6$ forward passes\footnote{except for the first iteration, which depends on the initial guess $\eta_0$ for line search.}.
    \item \looseness -1
    We then examine the relationship between critical sharpness and Hessian sharpness. Under the quadratic loss approximation, we show that critical sharpness can be written as a weighted sum of Hessian eigenvalues and the two measures coincide when the gradient aligns with the largest eigenvector of the Hessian. We also generalize this result to adaptive optimizers.
    
    \item We demonstrate that critical sharpness reliably captures well-documented Hessian sharpness phenomena, such as \emph{progressive sharpening} and the \emph{Edge of Stability}. This makes critical sharpness a practical and computationally efficient proxy for sharpness. Leveraging OLMo-2 checkpoints~\citep{walsh2025} across pre-training and mid-training, we demonstrate progressive sharpening persists at realistic scales throughout training, including models with up to 7 billion parameters.

    \item We introduce a relative measure of critical sharpness $\lambda_c^{1 \to 2}$ to quantify the curvature of the pre-training loss landscape during mid-training. By varying the mix of pre-training and math data, we show that increasing the proportion of pre-training data lowers relative critical sharpness, helping the model stay near the ``pre-trained basin.'' We further demonstrate that downstream performance depends on basin retention: GSM8K (math) benefits from leaving the basin ($\eta > 2 / \lambda_c^{\text{rel}}$), while MMLU (generic reasoning) performs better when the basin is retained ($\eta < 2 / \lambda_c^{\text{rel}}$). This enables us to prescribe data mixing strategies to balance performance on math and generic reasoning tasks, depending on the desired outcome.

    % \item Using a relative variant of critical sharpness $\lambda_c^{\text{rel}}$, we investigate the effect of pre-training dataset composition during mid-training. We find that increasing the proportion of pre-training data in the mid-training mix consistently lowers relative critical sharpness of the pre-training landscape, which in turn enables stable training at higher learning rates without losing what we have learnt in the previous pre-training stage. We also find a dichotomy: datasets in the mix that align with the 

\end{itemize}

    
\section{Characterizing Critical Sharpness: Theory and Empirics}

In this section, we characterize the dynamics critical sharpness both theoretically and empirically, by comparing its behavior with Hessian sharpness in controlled settings. 

\subsection{Setup}

Consider a model with trainable parameters $\bm{\theta} \in \mathbb{R}^n$, and let $L(\bm{\theta})$ denote the loss function with gradient $g(\bm{\theta}):=\nabla_\theta L(\bm{\theta})$ and Hessian $H(\bm{\theta}):=\nabla_\theta^2 L(\bm{\theta})$. The Hessian is a symmetric matrix with eigenvalues $\{\lambda_{i}^H\}_{i = 1}^n$ and corresponding eigenvectors $\{{u}_{i}(\bm{\theta})\}_{i = 1}^n$. The largest eigenvalue $\lambda_{\max}^H = {\max}_{i} \{\lambda_{i}^H \}$, referred to as \emph{Hessian sharpness}, quantifies the worst-case local curvature of the loss landscape. For adaptive optimizers, such as Adam, with pre-conditioner $P(\bm{\theta})$, the dynamics is governed by the pre-conditioned Hessian $P(\bm{\theta})^{-1/2}H(\bm{\theta})P(\bm{\theta})^{-1/2}$ with eigenvalues $\{\lambda^{PH}_i \}_{i=1}^n$ and eigenvectors $\{{v}_i(\bm{\theta})\}_{i=1}^n$. Note that both the eigenvalues and eigenvectors evolve as training progresses. For notational simplicity, we may occasionally omit the explicit $\bm{\theta}$ dependence when it is clear from the context.

\begin{figure}[!t]
    \centering
    \includegraphics[width=0.3\linewidth]{figures/landscape-edited.png}
    \caption{\emph{Comparison of different sharpness measures on an illustrative landscape featuring a sharp valley direction and a flat river direction.} Hessian sharpness quantifies the curvature along the sharpest direction of the landscape (the valley), while critical sharpness measures the curvature along the update direction $\Delta \bm{\theta}$.}
    \label{fig:landscape}
\end{figure}
\subsection{Critical Sharpness: A Scalable Measure of Curvature}

\begin{definition}[Critical learning rate $\eta_c$ and Sharpness $\lambda_c$]
\label{def:critical-lr-sharpness}   
Given an update direction $\Delta \bm{\theta}$, we define \emph{critical learning rate} as the smallest learning rate that causes the loss to increase in the next training step:
\begin{align}
    \eta_c = \min_{\eta > 0} \{ \eta \mid L(\bm{\theta} - \eta \Delta \bm{\theta} ) >  L(\bm{\theta})  \}.
\end{align}

Correspondingly, we define \emph{critical sharpness} as the scaled reciprocal of the critical learning rate $\lambda_c = 2 / \eta_c$.
\end{definition}

\paragraph{\textbf{Geometric Interpretation}:} 
Critical sharpness provides an intuitive, optimizer-aware measure of the local curvature of the loss landscape. Geometrically, it quantifies the ``natural length-scale'' of the landscape by estimating how far one can move in the current update direction without leaving the local region. \Cref{fig:landscape} compares critical sharpness with Hessian sharpness on an illustrative landscape.

% Critical sharpness quantifies the natural scale of the landscape by estimating the largest step size that can be taken along the update direction $\Delta \bm{\theta}$ without increasing the loss. 




\paragraph{\textbf{Efficient estimation of critical learning rate:}} Given the update direction $\Delta \bm{\theta}$ from training, we compute the critical learning rate $\eta_c$ using a two-phase line search procedure:
\begin{enumerate}
    \item \emph{Exponential Search:} Starting from an initial guess $\eta_0$, we iteratively increase (decrease) $\eta$ until the loss $L(\bm{\theta} - \eta \Delta \bm{\theta})$ exceeds (or falls below) the current loss $L(\bm{\theta})$. This quickly identifies an interval $[\eta_{\text{lower}}, \eta_{\text{upper}}]$ containing $\eta_c$.
    \item \emph{Binary Search:} We then refine this interval using a binary search. At each iterate, we evaluate the loss at the midpoint $\eta_{\text{mid}} = \frac{1}{2}(\eta_{\text{lower}} + \eta_{\text{upper}})$ and update the interval accordingly. This process is continued until the relative error $|1 - \eta_{\text{lower}}/ \eta_{\text{upper}}|$ falls below a specified threshold $\epsilon$. For $\epsilon = \frac{1}{2^k}$, the binary search converges in $k$ steps. 
\end{enumerate}

Finally, we approximate the critical learning rate using the mean of the final range $\eta_c \approx \frac{1}{2}(\eta_{\text{lower}} + \eta_{\text{upper}})$.
In practice, the exponential search typically requires only $1$–$2$ iterations (except for the first use, which depends on the initial guess $\eta_0$), and setting $\epsilon = 1/16$ (i.e., $k=4$) provides a reliable and efficient estimate of the critical learning rate in approximately $5$–$6$ forward passes (see \Cref{fig:sharpness_dynamics_cifar10_sgd}). This approach is highly scalable, as it relies solely on forward passes to evaluate $L(\bm{\theta} - \eta \Delta \bm{\theta})$ and leverages the same computational primitives as standard first-order distributed training. Notably, the line search scales linearly with the number of parameters $n$, in contrast to the quadratic scaling $\Theta(n^2)$ required for Hessian-based methods, and it does not suffer from the non-convergence issues associated with eigenvalue solvers at scale. \Cref{table:computational_complexity_sharpness}
In \Cref{appendix:critical-lr-estimation}, we provide the detailed algorithm and further discuss the design choices.


% In practice, we find that the exponential search takes $\sim 1-2$ iterations, except for the first iteration, which depends on the initial guess $\eta_0$ and we find that setting $\epsilon = 1/16$ (i.e., $k=4$) provides a reliable and efficient estimate of the critical learning rate (see \Cref{fig:sharpness_dynamics_cifar10_sgd}).
% This procedure reliably estimates critical learning rate in only $\sim 5-6$ fo
% This procedure is highly scalable, as it only required forward passes to evaluate $L(\bm{\theta} - \eta \Delta \bm{\theta})$ and utilizes the same computational primitives used in standard first-order distributed training. 
% Furthermore, this simple line search scales linearly with the number of parameters $n$, as compared to the quadratic scaling $\Theta(n^2)$ required for computing Hessian-vector-products and does not suffer from the non-convergence of eigenvalue solvers. In \Cref{appendix:critical-lr-estimation}, we provide the detailed algorithm and further discuss the design choices.

% This is much cheaper than sharpness estimation via eigenvalue solvers, which require many expensive Hessian-vector product computations and can suffer from convergence issues, whereas our method uses standard training operations and is straightforward to implement.




\subsection{The Relationship between Critical Sharpness and Hessian Sharpness}

\begin{figure*}[t]
     \centering
     \begin{subfigure}[b]{0.32\textwidth}
         \centering
         \includegraphics[width = \linewidth]{figures/cifar/sgd/cifar10_sharpness_xent_bs50000_lr3e-02.pdf}
         \caption{Batch size = 50,000 (GD)}
     \end{subfigure}
     \hfill
     \begin{subfigure}[b]{0.32\textwidth}
         \centering
         \includegraphics[width = \linewidth]{figures/cifar/sgd/cifar10_sharpness_xent_bs5000_lr3e-02.pdf}
         \caption{Batch size = 5000}
     \end{subfigure}
     \hfill
     \begin{subfigure}[b]{0.32\textwidth}
         \centering
         \includegraphics[width = \linewidth]{figures/cifar/sgd/cifar10_sharpness_xent_bs500_lr3e-02.pdf}
        \caption{Batch size = 500}
     \end{subfigure}
     \caption{Comparison of different sharpness measures for MLPs trained on CIFAR-10 image classification task using SGD with learning rate $\eta$. 
     Both critical and directional sharpness exhibit progressive sharpening and Edge of Stability, albeit some deviations from Hessian sharpness.
     The dashed line denotes the Edge of Stability threshold, given by $2/\eta$.}
     \label{fig:sharpness_dynamics_cifar10_sgd}
\end{figure*}


\looseness -1
Before interpreting the empirical behavior of critical sharpness, it is important to establish its theoretical relationship with Hessian sharpness. This connection provides a principled foundation for understanding what critical sharpness captures.
To this end, we consider the quadratic approximation of the loss function along $\Delta \bm{\theta}$:
\begin{align}
    L(\bm{\theta} - \eta \Delta \bm{\theta}) \approx L(\bm{\theta}) - \eta \Delta \bm{\theta}^T g(\bm{\theta}) + \frac{1}{2} \eta^2 \Delta \bm{\theta}^T H(\bm{\theta}) \Delta \bm{\theta}. \nonumber
\end{align}
Within this approximation, the loss will increase if the learning rate $\eta$ exceeds $2 / \lambda_{\text{dir}}$, where $\lambda_{\text{dir}}$ is the \emph{directional sharpness}~\citep{pan2022toward,roulet2024stepping}:
\begin{definition}[Directional Sharpness $\lambda_{\text{dir}}$]
The directional sharpness along the update direction $\Delta \bm{\theta}$ is:
\begin{align}
    \lambda_{\text{dir}} = \frac{\Delta \bm{\theta}^T H(\bm{\theta}) \Delta \bm{\theta} }{\Delta \bm{\theta}^T g(\bm{\theta})}.
\end{align}
\end{definition}

Directional sharpness $\lambda_{\text{dir}}$ serves as an analytically tractable approximation to the empirically measured critical sharpness. \Cref{fig:sharpness_dynamics_cifar10_sgd} shows that directional sharpness closely tracks critical sharpness throughout training. 
This alignment suggests that the local quadratic approximation well captures the complex dynamics in this simplistic setting.
For Gradient Descent (GD), the directional sharpness can be expressed as a weighted sum of Hessian eigenvalues~\citep{roulet2024stepping}:

% While critical sharpness captures the exact learning rate above which the loss increases through empirical observation, directional sharpness provides an analytical approximation based on the quadratic approximation. 


% \Cref{fig:sharpness_dynamics_cifar10_sgd} shows that critical sharpness and directional sharpness both exhibit progressive sharpening and EoS phenomena, and the two measures closely track each other. This alignment suggests that the local quadratic approximation well captures the complex dynamics of neural network optimization.
% Next, we establish a relationship between directional sharpness and Hessian sharpness.



\begin{result}[Relationship between Directional and Hessian Sharpness for Gradient Descent]
\label{prop:dir_hessian_rel_gd}
\looseness -1
For Gradient Descent (GD), the directional sharpness $\lambda_{\text{dir}}$ can be expressed as a weighted sum of the Hessian eigenvalues $\{\lambda_i^H\}_{i=1}^n$, where the weights quantify the alignment of the gradient with Hessian eigendirections $\{u_i\}_{i=1}^n$.
\begin{align}
    \lambda_{\text{dir}} = \frac{\sum_{i = 1}^n c_i^2 \lambda_i^H}{\sum_{i = 1}^n c_i^2} \leq \lambda_{\max}^H
\end{align}
where $c_i = \bm{g}^T \bm{u}_i$ is the projection of the gradient onto the $i^{th}$ eigenvector.
\end{result}

It follows that if the gradient $\bm{g}$ is perfectly aligned with the largest eigenvector, i.e., $ \bm{g} \propto \bm{u}_{\max}$, then directional sharpness coincides with Hessian sharpness. More generally, directional sharpness is always smaller than the Hessian sharpness. We now generalize this result to adaptive optimizers.

\begin{result}[Relationship between Directional and Hessian Sharpness for Adaptive optimizers]
\label{prop:dir_hessian_rel_adam}
For adaptive optimizers (such as Adam) with pre-conditioner $P$, the directional sharpness $\lambda_{\text{dir}}$ can be expressed as a weighted sum of the pre-conditioned Hessian eigenvalues $\{\lambda^{PH}_i\}_{i=1}^n$, where the weights quantify the alignment of the pre-conditioned gradient $P^{-1/2} \bm{g}$ with pre-conditioned Hessian eigendirections $\{\bm{v}_i\}_{i=1}^n$.
\begin{align}
    \lambda_{\text{dir}} = \frac{\sum_{i = 1}^n c_i^2 \lambda_i^{PH}}{\sum_{i = 1}^n c_i^2} \leq \lambda_{\max}^{PH}
\end{align}
where $c_i = P^{-1/2}\bm{g}^T \bm{v}_i$ is the projection of the gradient $g(\bm{\theta})$ onto the $i^{th}$ eigenvector $\bm{v}_i$ of the pre-conditioned Hessian.
\end{result}
% For vanilla GD $\Delta \bm{\theta} = g(\bm{\theta})$, if the gradient aligns with the largest eigenvector of the Hessian, i.e., 
% $g(\bm{\theta}) = \|g(\bm{\theta})\| \bm{u}_{\max}$, then the directional sharpness equals the Hessian sharpness. In contrast, these two measures deviate when the gradient has significant contribution along other eigendirections.
% %
% As a consequence, the gap between Hessian and directional sharpness tells us about the alignment of gradient with the eigenspace.
% %
% The above result generalizes to adaptive optimizers \textemdash the directional sharpness is equal to the pre-conditioned sharpness $\lambda^{PH}_{\max}$ when the pre-conditioned gradient $P^{-1/2} g(\bm{\theta})$ aligns with the largest eigenvector of the pre-conditioned Hessian.
% For vanilla gradient descent, it relates the Hessian sharpness, while for adaptive optimizers, it captures the pre-conditioned sharpness dynamics. As a result, it is an optimizer-aware measure of the local curvature of the loss landscape. 
% \paragraph{\textbf{Empirical comparison of the different sharpness measures:}} 
% Explain the difference between different sharpness

\looseness -1
Together, these results establish a connection between directional sharpness and Hessian sharpness, and indicate that the two measures diverge when the alignment between the gradient and the top eigendirection is small. In turn, directional and critical sharpness coincide when the loss surface is approximately locally quadratic and the gradient is primarily along the top eigendirection. 
We now examine the differences and similarities between these three sharpness measures empirically.
\Cref{fig:sharpness_dynamics_cifar10_sgd} compares Hessian sharpness, directional sharpness, and critical sharpness of FCNs trained on the CIFAR image classification task using SGD with fixed learning rate $\eta$ and different batch sizes.
In the full-batch regime, Hessian sharpness exhibits progressive sharpening, eventually reaching the Edge of Stability (EoS) threshold. In contrast, both directional sharpness and critical sharpness remain nearly constant during the early stages of training, followed by a abrupt rise to the EoS threshold.
At smaller batch sizes, however, all three sharpness measures display a gradual increase from the onset of training. Notably, after crossing the EoS threshold, Hessian sharpness tends to oscillate above the threshold, whereas both critical sharpness and directional sharpness oscillate more closely around it.
Overall, critical sharpness exhibits progressive sharpening and Edge of Stability, albeit some differences compared to Hessian sharpness. In the following section, we extend our analysis to more realistic, large-scale training settings.



% The gap between sharpness and critical sharpness describes the difference between eta_c and sharpness
% Our work generalizes Kalra and Barkeshli's result that critical learning rate >> 2/sharpness
% 1. all sharpness metrics exhibit progressive sharpening and EoS
% 2. In the full batch setting, the increase in directional and critical sharpness is abrupt during, while sharpness increases slowly. Interpretation
% 3. The rate of progressive sharpening reduces with batch size as in Pennington et al.

\begin{table}[!t]
    \centering
    \normalsize
    \renewcommand{\arraystretch}{1.4}
    \begin{tabular}{l|c}
        \hline
        \textit{Sharpness Measure} & \textit{Computational Complexity} \\
        \hline
        {Hessian sharpness} ($\lambda^H_{\max}$) & $\Theta(p \cdot n^2)$ \\
        \hline
        {Directional sharpness} ($\lambda_{\text{dir}}$) & $\Theta(n^2)$ \\
        \hline
        {Critical sharpness} ($\lambda_c$) & $\mathcal{O}(m \cdot n)$ \\
        \hline
    \end{tabular}
    \caption{\emph{Computational complexity of different sharpness measures.} Hessian sharpness requires $\Theta(p \cdot n^2)$ operations, where $p$ is the number of eigenvalue solver iterations.  Directional sharpness requires $\Theta(n^2)$ operations. In contrast, critical sharpness can be computed in $\Theta(m \cdot n)$ operations, where $m$ is the number of forward passes, making it more scalable to large models.}
    \label{table:computational_complexity_sharpness}
\end{table}


\section{Critical Sharpness Dynamics at Scale}



% Having studied the fundamental properties of critical sharpness in simplistic settings, we now extend our analysis to the realistic setting of LLM training. 
Modern large-scale models are typically trained using Adam with weight decay~\citep{adamwloshchilov2018}, which helps mitigate training instabilities~\citep{dangelo2024why}. To analyze sharpness dynamics in this context, we first analyze the stability threshold for common optimizers with weight decay:
\begin{result}[Stability threshold for optimizers with weight decay]
For gradient descent and adaptive optimizers such as Adam, adding weight decay shifts the EoS threshold by the decay strength $\gamma$:
\begin{align}
    \text{For GD:} \quad & \lambda^H_{\max} = \frac{2}{\eta} - \gamma \nonumber \\
    \text{For Adam:} \quad & \lambda^{PH}_{\max} = \frac{2 + 2\beta_1}{(1 - \beta_1)\eta} - \gamma
\end{align}
where $\eta$ is the learning rate, $\beta$ and $\beta_1$ are momentum parameters, and $\gamma$ is the weight decay strength.
\end{result}

\looseness -1
We now analyze GPT-style Transformers trained on language pre-training task on FineWebEdu using AdamW with Warmup-Stable-Decay (WSD) schedule~\citep{hu2024minicpm}.
\Cref{fig:sharpness_dynamics_nanogpt_pretraining} compares the dynamics of critical and pre-conditioned sharpness for three different learning rates. The pre-conditioned sharpness $\lambda^{PH}_{\max}$ exhibits progressive sharpening while following the learning rate schedule throughout training \textemdash it is pushed down during the warmup phase, stays constant at the stability threshold during the stable phase, and increases again when the learning rate is decayed. Critical sharpness and directional sharpness follow similar trends throughout training \textemdash they exhibit clear progressive sharpening and EoS behavior, while oscillating below the EoS threshold during the stable phase and capture the increase in pre-conditioned sharpness during the decay phase.

\begin{figure*}[t]
     \centering
     \begin{subfigure}[b]{0.32\textwidth}
         \centering
         \includegraphics[width = \linewidth]{figures/pre-training/sharpness/presharp_fineweb_gpt_d12_h12_n768_AdamW_Tw2000_r1.0_Ts6000_cosine_p1.0_T10000_ga64_lrinf_wd0.0_bs16_b0.9_b0.95_eps1e-08_gc0.0.pdf}
         \caption{}
     \end{subfigure}
     \begin{subfigure}[b]{0.32\textwidth}
         \centering
         \includegraphics[width = \linewidth]{figures/pre-training/sharpness/dirsharp_fineweb_gpt_d12_h12_n768_AdamW_Tw2000_r1.0_Ts6000_cosine_p1.0_T10000_ga64_lrinf_wd0.0_bs16_b0.9_b0.95_eps1e-08_gc0.0.pdf}
         \caption{}
     \end{subfigure}
     \begin{subfigure}[b]{0.32\textwidth}
         \centering
         \includegraphics[width = \linewidth]{figures/pre-training/sharpness/critsharp_fineweb_gpt_d12_h12_n768_AdamW_Tw2000_r1.0_Ts6000_cosine_p1.0_T10000_ga64_lrinf_wd0.0_bs16_b0.9_b0.95_eps1e-08_gc0.0.pdf}
         \caption{}
     \end{subfigure}
     \caption{Dynamics of pre-conditioned, directional, and critical sharpness during GPT-style Transformer training on FineWebEdu with AdamW. Critical sharpness tracks pre-conditioned sharpness throughout training, making it an effective proxy.}
    \label{fig:sharpness_dynamics_nanogpt_pretraining}
\end{figure*}



\looseness -1
While progressive sharpening has been consistently documented at small scales, especially in full-batch settings, its manifestation and relevance at large scales, particularly in the context of online LLM training, remains largely unexplored. Having established that critical sharpness serves as an efficient and reliable proxy for Hessian sharpness phenomena, we are now in the position to investigate whether progressive sharpening persists in large-scale LLM training.
%
To this end, we analyze the publicly available OLMo-2 $7$B checkpoints~\citep{walsh2025}, throughout training. 

\looseness -1
The OLMo-2 models are trained using a two-stage curriculum approach. During the pre-training stage ($>90\%$ of compute), the model is trained for $4$T tokens on a web-sourced corpus. This pre-training data primarily consists of the DCLM dataset~\citep{li2025datacomplmsearchgenerationtraining},  a diverse web-sourced corpus, supplemented by smaller proportions of additional high-quality documents.
Following pre-training, the model undergoes a pre-training stage, where it is trained on a curated mixture of datasets referred to as the Dolmino mix for $50$B tokens.
This mix includes academic papers, Wikipedia articles, instruction-following data, StackExchange documents, and approximately $50\%$ of the original pre-training dataset (DCLM).
%
During the pre-training stage, the model is trained using a learning rate schedule consisting of $2,000$ steps of warmup, followed by a cosine decay down to one-tenth of its peak value\footnote{For the OLMo $7$B models, the cosine decay schedule is designed to reach one-tenth of the peak learning rate at $5$T tokens, but is truncated at $4$T tokens.}. By comparison, during the mid-training stage, the learning rate is linearly decayed to zero. As the learning rate is continuously decreased throughout training after the warmup stage, we expect sharpness to progressively increase as training proceeds.
To confirm this, we examine how the critical sharpness evolves across pre-training and mid-training using the publicly available checkpoints (see \Cref{appendix:details-olmo-checkpoints} for the details).
\Cref{fig:crit_sharp_olmo_checkpoints} shows that critical sharpness decreases rapidly during the initial phase of training, but then continually increases (progressive sharpening) throughout both the pre-training and mid-training stages. This result provides the first empirical evidence of progressive sharpening at scale in practical LLM training settings.


\begin{figure*}[t]
     \centering
     \begin{subfigure}[b]{0.32\textwidth}
         \centering
         \includegraphics[width = \linewidth]{figures/417-420-curvature-scale/critical_sharpness_vs_dclm_ratio_exp420.pdf}
         \caption{}
     \end{subfigure} 
     \begin{subfigure}[b]{0.32\textwidth}
         \centering
         \includegraphics[width = \linewidth]{figures/417-420-curvature-scale/critical_sharpness_vs_dclm_ratio_exp417.pdf}
         \caption{}
     \end{subfigure}
     \caption{Critical Sharpness of OLMo $7$B exhibits progressive sharpening throughout pre-training and mid-training. In the mid-training figure, the band around the mean trend shows the deviation across the three runs. }
    \label{fig:crit_sharp_olmo_checkpoints}
\end{figure*}



\section{How much Pre-training data is needed to avoid Catastrophic forgetting?}

During finetuning, neural networks are prone to \emph{catastrophic forgetting}, where the model performance degrades on the pretraining dataset and benchmarks as the model adapts to the new task~\citep{MCCLOSKEY1989109,lou2023_forgetting,mcleish2025teachingpretrainedlanguagemodels}. To mitigate this, several strategies have been proposed~\citep{Lange2022}, with mixing samples from the pretraining data being the most effective~\citep{ROBINS01061995}. This practice is reflected in large-scale training, such as~\cite{walsh2025}, where the mid-training data consists of approximately $50\%$ pre-training data. However, it remains unclear what fraction of pre-training data is sufficient to effectively prevent catastrophic forgetting.
To this end, we leverage critical sharpness to systematically examine the effect of adding pretraining data during OLMo mid-training and provide actionable guidance for selecting the pre-training fraction without exhaustive grid search.


The goal of mid-training or finetuning is to improve the performance on specialized domains or tasks while preserving generic capabilities acquired during pretraining. Intuitively, this requires the model to adapt to the new task, while staying within the ``pretraining basin'' \textemdash a region of the parameter space where the pre-training loss remains low.  Leaving this basin would be marked by an increase in the pre-training loss, which naturally relates to critical sharpness. 


To formalize this intuition, we define relative critical learning rate and sharpness, as follows:
\begin{definition}[Relative Critical learning rate $\eta_c^{1 \to 2}$ and Sharpness $\lambda_c^{1 \to 2}$]
\label{def:relative-critical-lr-sharpness}   
Consider a model evaluated on Task 1 with loss $L_1(\bm{\theta})$, while being trained on Task 2 with loss $L_2(\bm{\theta})$ and update direction $\Delta \bm{\theta}_2$. Then, the \emph{relative critical learning rate} is the smallest learning rate for which taking a step in the direction $\Delta \bm{\theta}_2$ increases the loss $L_1(\bm{\theta})$ in the next training step:
\begin{align}
    \eta_c^{1\to2} = \min_{\eta > 0} \{ \eta \mid L_1(\bm{\theta} - \eta \Delta \bm{\theta}_2 ) >  L_1(\bm{\theta})  \}.
\end{align}
The corresponding \emph{relative critical sharpness} is $\lambda_c^{1\to2} = 2 / \eta_c^{1\to2}$.
Here, Task 1 can be pre-training a LLM on a general text corpus, and Task 2 can be finetuning the pre-trained model on a specialized dataset, such as math. 
 
\end{definition}

To assess the impact of pre-training data fraction in the training mix, we consider the OLMo-2 $7$B pre-trained checkpoint as our starting point, and further trained it on a mixture (Task 2) composed of DCLM (pre-training data) and the math subset of Dolmino mix~\citep{walsh2025}. We then examine how the relative critical sharpness of the pre-trained checkpoint varies with the pre-training mix fraction\footnote{When measuring the relative critical sharpness, we do not update the model parameters.}. 
\Cref{fig:pretrain_mix}(a) shows how the relative critical sharpness varies with the DCLM ratio in the training mix for several evaluation tasks. When the mix contains mostly math data (low DCLM ratio), the sharpness for DCLM is an order of magnitude higher than for other tasks, indicating that the pre-training loss landscape is particularly sharp and sensitive to updates in this regime. As more pre-training data is added, the relative critical sharpness for most tasks decreases, suggesting that the landscapes align. However, when the DCLM ratio approaches one, the sharpness for downstream tasks such as Math, GSM8K, and MMLU increases, meaning that these tasks become the limiting factor for the maximum stable learning rate. Notably, there is an intermediate DCLM ratio (around 0.7) where the sharpness curves for different tasks intersect, representing a sweet spot that allows for the largest possible learning rate without being constrained by any single task. This result highlights the importance of balancing pre-training and downstream data to optimize training stability and efficiency.



Next, we evaluate the impact of the pre-training mix by training the checkpoint for 1B tokens and measuring downstream accuracy on GSM8K and MMLU benchmarks, as shown in \Cref{fig:pretrain_mix}(b, c). In these plots, red indicates a decrease in performance relative to the checkpoint, white indicates no change, and blue indicates an improvement. 

We observe two distinct post-training regimes: \emph{in-basin} and \emph{out-of-basin}. When post-training in-basin (learning rate below the critical learning rate), there is no performance deterioration across both benchmarks

We observe a natural trade-off: GSM8K accuracy is typically maximized outside the pre-training basin (i.e., with low DCLM ratio and high learning rate), but this comes at the expense of MMLU performance, which is best preserved within the pre-training basin. Therefore, if the sole objective is to improve on GSM8K, then training primarily on Math with a large learning rate is effective, though it may lead to forgetting on other tasks, such as MMLU. In contrast, if the goal is to improve on GSM8K while maintaining MMLU performance, it is essential to include DCLM in the mix, with a sweet spot emerging around a DCLM ratio of $0.6$ and a learning rate of 3e-5. Remarkably, this is quite close to the optimum suggested by the critical sharpness analysis. This result highlights the importance of balancing task-specific and general data in the training mix \textemdash once the loss on a task increases and its data is absent from the mix, performance on that task can degrade significantly due to unconstrained updates along the gradient direction.


\section{Related Works}


% \paragraph{\textbf{Mode Connectivity:} }\cite{Garipov18_loss_surfaces} showed that the minima of neural network are connected by simple, yet non-linear paths. Interestingly, after early training, the resulting minima, obtained under different samples of SGD noise, can become linearly connected~\citep{pmlr-v119-frankle20a}.
    
In neural network training, if the learning rate exceeds the threshold $\sim 2/\lambda^H_{\max}$
max, it causes the loss to increase in the next training step~\cite{Wu2018_how_sgd}. However, \cite{kalra2023phase} demonstrated that the empirical critical learning rate can be much higher than this theoretical threshold derived from convex analysis, with critical sharpness reaching up to $40/\lambda^H_{\max}$ when sharpness decreases significantly during training.
Our work extends the work by \cite{kalra2023phase} showing that a similarly large critical learning rate can also arise in cases when sharpness increases during training, due to contributions from other eigendirections (\Cref{fig:sharpness_dynamics_cifar10_sgd}). In addition to training instabilities, Hessian sharpness exhibits progressive sharpening during training. In particular, \cite{cohen2021gradient} showed that sharpness increases when the learning rate is decayed, which we corroborate with additional analysis at unprecedented scales (\Cref{fig:sharpness_dynamics_nanogpt_pretraining,fig:crit_sharp_olmo_checkpoints}).

\begin{figure*}[t]
     \centering
     \begin{subfigure}[b]{0.31\textwidth}
         \centering
         \includegraphics[width = \linewidth]{figures/440-pre-train-mix-init/critical_sharpness_vs_dclm_ratio_exp440.pdf}
         \caption{}
     \end{subfigure}
     \begin{subfigure}[b]{0.32\textwidth}
         \centering
         \includegraphics[width = \linewidth]{figures/433-pre-train-mix-training/gsm8k_accuracy_landscape_exp433_stage1.pdf}
         \caption{}
     \end{subfigure}
     \begin{subfigure}[b]{0.32\textwidth}
         \centering
         \includegraphics[width = \linewidth]{figures/433-pre-train-mix-training/mmlu_accuracy_landscape_exp433_stage1.pdf}
         \caption{}
     \end{subfigure}
     \caption{(a) Relative critical sharpness (\Cref{def:relative-critical-lr-sharpness}) for various evaluation tasks. The shaded region around the mean trends denotes the variation across batches. (b, c) GSM8K and MMLU accuracy as a function of pre-training (DCLM) mix ratio and learning rate. Red indicates a decrease in performance relative to the checkpoint, white indicates no change, and blue indicates an improvement. The black dashed line denotes the smallest critical learning rate among at the pre-trained checkpoint across tasks from (a) and white dashed line denotes the critical learning rate for the corresponding task.}
    \label{fig:pretrain_mix}
\end{figure*}


The works most closely related to ours are those by \cite{kalra2024why} and \cite{roulet2024stepping}. \cite{kalra2024why} use critical learning rate to set the initial learning rate during warmup. By comparison, \cite{roulet2024stepping} propose setting the learning rate at `edge' throughout training, i.e., $\eta_t = 2/\lambda_{\text{dir}}$) and show that it can outperform constant learning rate training, but not typical learning rate schedules, involving learning rate warmup followed by decay. ~\citep{Vaswani2019Painless} also study setting the learning rate as $\eta = 1 / \lambda_{\text{dir}}$. In contrast, we use the critical learning rate to study the curvature dynamics of large scale models and examine the effect of pre-training data on catastrophic forgetting.

Catastrophic forgetting is a long standing problem in neural networks, first identified by~\cite{MCCLOSKEY1989109}, where model performance degrades on previously learned task as it adapts to the new task. Approaches to mitigate catastrophic forgetting can be categorized into (i) regularization methods~\citep{Kirkpatrick2017,Ahn2019}, (ii) ensembling and parameter isolation~\cite{rusu2016progressive,Aljundi2017}, and  (iii) rehearsal~\citep{ROBINS01061995,Lopez-Paz2017}. Among these, rehearsal, mixing samples from the pretraining data, has become the most widely adopted strategy due to its simplicity and effectiveness. Recent works have demonstrated that catastrophic forgetting persists in LLM finetuning~
\citep{lou2023_forgetting,huang-etal-2024-mitigating,scialom-etal-2022-fine}. Of particular similarity to our work is by ~\cite{chen2025understandingpretrainingfinetuningloss} who show that if finetuning retains previously learned capabilities, it remains in `most-case' and `worst-case' basins.

    
\section{Discussion and Conclusion}


In this work, we introduced and analyzed critical sharpness, a scalable and practical measure for analyzing the training dynamics of LLMs at unprecedented scale. Our results demonstrate that critical sharpness reliably captures key Hessian sharpness phenomena, such as progressive sharpening and the Edge of Stability, while being far more computationally efficient. By leveraging this measure, we provided the first empirical evidence of progressive sharpening in models with up to 7B parameters in practical LLM training settings, spanning both pre-training and mid-training regimes.

Another central contribution of our study is the introduction of relative critical sharpness, which quantifies the curvature of the pre-training loss landscape during mid-training or fine-tuning. Through systematic experiments on OLMo-2 checkpoints, we showed that increasing the fraction of pre-training data in the training mix consistently lowers relative critical sharpness, thereby stabilizing training and facilitating training at larger learning rates without catastrophic forgetting. Our analysis revealed a natural trade-off: while specializing on a downstream task (e.g., Math) with a large learning rate and little pre-training data can maximize performance on that task, it often leads to rapid degradation on more general benchmarks (e.g., MMLU). In contrast, including sufficient pre-training data in the mix preserves general capabilities, and we identified a ``sweet spot'' in the data mix that optimally balances specialization and retention, as predicted by the critical sharpness analysis.
Our findings have practical implications for the designing data mix strategies and learning rate schedules in large-scale model training. 

Critical sharpness enables practitioners to make informed decisions about training stability, data composition, and the risk of catastrophic forgetting, without the prohibitive cost of Hessian analysis. More broadly, our results highlight the importance of understanding and monitoring curvature dynamics for developing robust and effective optimization strategies in modern neural network training.

Looking forward, critical sharpness opens several avenues for future research. It can be extended to analyze training instabilities in reinforcement learning objectives and structured quantization, or to guide adaptive learning rate schedules and curriculum design. Additionally, its optimizer-aware nature makes it a promising diagnostic took for emerging training paradigms and architectures. We hope that this work will inspire further exploration of scalable curvature measures and their role in large-scale neural network training.

% \paragraph{Negative results:} Sharpness does not tell us about the optimal decay schedule.

\section*{Acknowledgements}

We would like to thank Tianyu He, Sean Mcleish, Konstantin Mishchenko, Aaron Dafezio, Maissam Barkeshli, Benjamin Therien, Jesse Dodge for helpful discussions. 
\bibliographystyle{assets/plainnat}
\bibliography{Meta-AI/ref}

\clearpage
\newpage
\beginappendix


\section{Estimating Critical Learning Rate Using Forward Passes}
\label{appendix:critical-lr-estimation}

This section provides further details on how to measure critical learning rate just using forward passes. We generalize the line search approach by \cite{kalra2024why}

\begin{algorithm}[H]
\caption{Exponential Search}
\label{alg:exponential_search}
\begin{algorithmic}[1]
\STATE \textbf{Input:} Update direction $\Delta \bm{\theta}$, Initial loss $L(\bm{\theta})$, initial guess $\eta_0$, maximum iterations $N_{\max}$
\STATE \textbf{Output:} Interval $[\eta_{\text{lower}}, \eta_{\text{upper}}]$ containing the critical learning rate $\eta_c$
\STATE $\eta \gets \eta_0$
\STATE $i \gets 1 $ \hfill // Iteration
\STATE $L(\bm{\theta} - \eta \Delta \bm{\theta}) \gets \text{ComputeLoss}(\eta, \Delta \bm{\theta})$
\IF{$L(\bm{\theta} - \eta \Delta \bm{\theta}) < L(\bm{\theta})$}
    \STATE $\text{dir} \gets +1$ \hfill // Increase learning rate
\ELSE
    \STATE $\text{dir} \gets -1$ \hfill // Decrease learning rate
\ENDIF
\WHILE{$i < N_{\max}$}
    \STATE $\eta \gets \eta \times 2^{\text{dir}}$
    \STATE $L(\bm{\theta} - \eta \Delta \bm{\theta}) \gets \text{ComputeLoss}(\eta, \Delta \bm{\theta})$
    \STATE $i \gets i + 1$
    \IF{$\text{dir} = +1$ \AND $L(\bm{\theta} - \eta \Delta \bm{\theta}) > L(\bm{\theta})$}
        \STATE $\eta_{\text{lower}} \gets \eta / 2$
        \STATE $\eta_{\text{upper}} \gets \eta$
        \STATE \textbf{return} $[\eta_{\text{lower}}, \eta_{\text{upper}}]$
    \ENDIF
    \IF{$\text{dir} = -1$ \AND $L(\bm{\theta} - \eta \Delta \bm{\theta}) < L(\bm{\theta})$}
        \STATE $\eta_{\text{lower}} \gets \eta$
        \STATE $\eta_{\text{upper}} \gets 2\eta$
        \STATE \textbf{return} $[\eta_{\text{lower}}, \eta_{\text{upper}}]$
    \ENDIF
\ENDWHILE
\STATE \textbf{return} $[\eta, \eta]$
\end{algorithmic}
\end{algorithm}

Given the update direction $\Delta \bm{\theta}$ from training, we compute the critical learning rate $\eta_c$ using a two-phase line search procedure:

\paragraph{\textbf{Exponential Search:}} Starting from an initial guess $\eta_0$, we iteratively double (half) the learning rate $\eta$ until the loss $L(\bm{\theta} - \eta \Delta \bm{\theta})$ exceeds (or falls below) the current loss $L(\bm{\theta})$. This quickly identifies an interval $[\eta_{\text{lower}}, \eta_{\text{upper}}]$ containing $\eta_c$, with $\eta_{\text{upper}} = 2 \eta_{\text{lower}}$. Notably, each iteration requires only a single forward pass to evaluate the loss. The full algorithm is provided in 

The number of exponential iterations depends on how close the initial guess $\eta_0$ is to the true critical learning rate $\eta_c$. After the first iteration, we update the initial guess to our current estimate of $\eta_c$, i.e., $\eta_0 = \eta_c$. In the subsequent steps, only $1-2$ exponential steps are needed, since $\eta_c$ tends to remain stable. However, if the training exhibits a large instability, causing the landscape to drastically change, more steps may be necessary.
To prevent indefinite iterations in such edge cases, we cap the exponential search to at most $40$ iterations, which corresponds to a increasing or decreasing the learning rate by a factor of $10^{12}$ relative to the initial guess. 

    
\textbf{Binary Search:} We then refine this interval using a binary search. At each iterate, we evaluate the loss at the midpoint $\eta_{\text{mid}} = \frac{1}{2}(\eta_{\text{lower}} + \eta_{\text{upper}})$ and update the interval accordingly. This process is continued until the relative error $|1 - \eta_{\text{lower}}/ \eta_{\text{upper}}|$ falls below a specified threshold $\epsilon$. For $\epsilon = \frac{1}{2^k}$, the binary search converges in $k$ steps. 

Finally, we approximate the critical learning rate using the mean of the final range $\eta_c \approx \frac{1}{2}(\eta_{\text{lower}} + \eta_{\text{upper}})$.
In practice, the exponential search typically requires only $1$–$2$ iterations (except for the first use, which depends on the initial guess $\eta_0$), and setting $\epsilon = 1/16$ (i.e., $k=4$) provides a reliable and efficient estimate of the critical learning rate in approximately $5$–$6$ forward pass. 

\section{Experimental Details}
\label{appendix:experimental_details}

This section provides details about the different experimental setups in the paper.

\subsection{CIFAR Image Classification}
\label{appendix:details-cifar-iamge-classification}

We considered Fully Connected Networks (FCNs) with four layers and width $512$, trained on the CIFAR-10 image classification task using SGD with a constant learning rate $\eta = 3e-02$ three different batch sizes $B \in \{500, 5,000, 50,000\}$. The batch size of $50,000$ corresponds to full batch gradient descent.

\subsection{GPT Pre-training}
\label{appendix:details-gpt-pretraining}

\subsection{Critical Sharpness analysis of OLMo checkpoints}
\label{appendix:details-olmo-checkpoints}

\section{The Relationship between Directional and Hessian sharpness}
\label{appendix:directional_hessian_sharpness}

In this section, we provide justification for \Cref{prop:dir_hessian_rel_gd,prop:dir_hessian_rel_adam}.




\end{document}